\chapter{実験設計}

\section{本報告で実験として扱う条件}

本報告は、次の問いの少なくとも一つに答える記録だけを実験として採用する。

\begin{enumerate}
  \item 本番運用モデルは実際に学習され、教示動作誤差が継続的に低下したか。
  \item データ、カメラ、計算資源、学習方式を何通り試し、どの実行が十分な段階まで進んだか。
  \item キャップなし、方向、反転、復帰、自由回転キャップなどの重要条件を網羅したか。
  \item モデルは3視野を実際に利用したか。
  \item 強化学習研究が、学習・保存・オフライン検証までの閉ループを形成したか。
\end{enumerate}

目的や比較条件がなく、単にファイルが存在するだけの試行は結果に含めない。プロセスの終了記録はモデルの有効性を意味せず、1件の教示フレームはロボットの成功を意味せず、学習損失の低下はタスク成功率を意味しない。

\section{学習試行の系譜}

2026年4月1日から5月31日までの学習基盤記録を、ボトル分別、開栓、洗浄、挿入に関する明示語で抽出した結果、41回の実行試行、33の異なる実行名、14設定を確認した。ボトル分別基礎モデルが23回、洗浄・挿入探索が18回である。\ev{E13}

\figfull[.96\textwidth]{experiment_funnel.pdf}{学習試行のファネル。ボトル分別基礎モデルと洗浄・挿入探索で合計41回を試した。これは工学的反復量と実行安定性を示す図であり、実機能力を示す図ではない。}

\begin{table}[H]
\centering
\caption{学習試行マトリクス}
\begin{tabularx}{\textwidth}{>{\bfseries}p{34mm}p{23mm}p{29mm}p{28mm}X}
\toprule
方向 & 試行数 & 異なる設定 & 1万ステップ到達 & 目的と限界\\
\midrule
ボトル分別基礎モデル & 23 & 10 & 9 & バッチ、全量・軽量学習、カメラ、データ構成を探索；学習時の乱数初期値は1種類\\
洗浄・挿入探索 & 18 & 4 & 3 & 洗浄と将来の挿入を探索；目的が異なり、損失を分別と直接比較しない\\
合計 & 41 & 14 & 12 & 中断33、失敗5、終了記録3；終了記録はタスク成功ではない\\
\bottomrule
\end{tabularx}
\end{table}

バッチサイズ4、8、32、128、256、512を試しており、メモリ、データ読み込み、分散学習、実行安定性の確立に大きな工数を投入したことが分かる。確認できた実行はすべて同じ乱数初期値42であり、複数条件の平均や信頼区間はまだ報告できない。

\section{本番運用モデルの学習}

本番学習は25件の一意な確定版データ、29件のサンプリング項目、バッチサイズ256、NVIDIA H200 GPU 4基を用い、6,059件の履歴を残した。記録されたのは学習損失、勾配ノルム、モデル規模の変化、データ待ち時間、1ステップ時間であり、検証損失、試験指標、標準手順に基づく実機結果は記録されていない。\ev{E08}

\begin{table}[H]
\centering
\caption{本番学習で利用可能な指標}
\begin{tabularx}{\textwidth}{>{\bfseries}p{38mm}p{24mm}X}
\toprule
指標 & 記録 & 答えられる問い\\
\midrule
学習時の動作生成損失 & あり & 教示動作へ近づいているか\\
勾配・パラメータノルム & あり & 数値的に安定しているか\\
学習時間・ステップ & あり & 学習工数と保存位置\\
検証損失 & なし & 学習・検証差は判断不能\\
試験成功率 & なし & 未見場面への一般化は判断不能\\
実機工程別完了率 & なし & 失敗段階は定量化不能\\
標準試験に基づく処理量 & なし & 約2本/分は現場概算に限定\\
\bottomrule
\end{tabularx}
\end{table}

\section{現場展示の観測}

現場展示では、把持、開栓、本体とキャップの分別投入までの一連の反復を観測した。観測者は連続稼働が1時間を超え、平均約2本/分で、ランダムな卓上位置と一定のボトル変化に初期的な適応を示したと概算している。\ev{E10}

完全な映像、逐次計数表、試験台帳、固定初期条件、失敗時刻がないため、これらは「現場観測による概算」として扱う。成功率、標準偏差、信頼区間は算出せず、注意記録の時間範囲とも混同しない。

\section{強化学習探索}

強化学習研究は別の洗浄データを使用した。835軌跡、354,835原フレームから、切出し後299,347フレーム、238,816学習区間を構成した。人手で付与した終端ラベルは、正例137件、負例698件である。\ev{E15} これは人手編集ラベルであり、自動物理成功判定でも、基礎分別モデルの実機成功率でもない。

研究ラウンド28--32は固定11本のオフライン検証軌跡を使用し、独立試験集合はない。最終ラウンドの動作最適化・価値推定保存モデルは61,541,926パラメータで、最適化状態と目標ネットワークを含む。\ev{E16}

\begin{table}[H]
\centering
\caption{強化学習探索のエビデンス成熟度}
\begin{tabularx}{\textwidth}{>{\bfseries}p{31mm}p{24mm}X}
\toprule
工程 & 状態 & 説明\\
\midrule
実データと人手評価 & 完了 & 835軌跡；評価品質は人手ラベルに依存\\
学習データ生成 & 完了 & 238,816学習区間\\
モデル学習・保存 & 完了 & 連続5ラウンド、保存物を読み取り可能\\
固定オフライン検証 & 完了 & 11軌跡；小規模で独立試験なし\\
同条件の基礎モデル比較 & 未完了 & 強化学習の優位性を判断できない\\
実機での増分確認 & 未完了 & 実機能力が向上したとは記述しない\\
\bottomrule
\end{tabularx}
\end{table}
